\section{Introduction}
\label{sec:introduction}

Quantum chaos is conventionally diagnosed in the spectrum of a Hamiltonian.
After resolving symmetries and unfolding the density of states, chaotic levels
repel and their long-range correlations approach a Wigner--Dyson random-matrix
ensemble \cite{bohigas1984,atas2013}.
The connected spectral form factor packages the same two-level correlations
into a ramp and plateau.
The eigenstate thermalization hypothesis (ETH) complements this spectral
picture by describing matrix elements of observables in the energy basis
\cite{deutsch1991,srednicki1994,dalessio2016}.
All three constructions use nonzero energy differences, so exact degeneracy
opens a complementary regime in which state-space geometry carries the
resolving information.  For \(D\) states sharing the energy \(E_0\),
\begin{equation}
 \begin{aligned}
 Z_E(t)&=\sum_{a=1}^{D}e^{-itE_a}=D e^{-itE_0},\\
 K_{E,\mathrm{raw}}(t)&=\frac{|Z_E(t)|^2}{D}=D ,
 \end{aligned}
 \label{eq:energy-silence}
\end{equation}
and the ensemble-connected form factor is identically zero.  These identities
exhaust the internal energy statistic while preserving the full projector and
its response over parameter space.

The alternative studied here is geometric.
An isolated degenerate eigenspace defines a vector bundle over Hamiltonian
parameter space.
Its adiabatic transport is governed by the Wilczek--Zee connection and the
matrix-valued Berry curvature \cite{kato1950,berry1984,wilczekzee1984};
the symmetric partner is the quantum metric
\cite{provostvallee1980,kolodrubetz2017}.
Adiabatic eigenstate deformations are already known to be sensitive probes of
integrability breaking \cite{pandey2020}, and the geometry of parameter space
can distinguish ergodic from integrable random-matrix regimes
\cite{sharipov2024}.
Most directly, Chen, Colin-Ellerin, Mamroud, and Papadodimas proposed
non-Abelian Berry curvature as an intrinsic chaos diagnostic for exactly
degenerate BPS microstates \cite{chen2026}.
Their contrast is between structured, often vanishing curvature for smooth
horizonless sectors and random-matrix-like curvature for supersymmetric black
hole microstates.

We test that organizing principle in a purely condensed-matter fractional
topological manifold generated by a local frustration-free parent.
The manifold is the kernel of a positive semidefinite bosonic
fractional-quantum-Hall parent Hamiltonian.
Its degeneracy is exact and algebraic, and its projector is separated from the
complement by a gap.
This setting isolates the proposed geometric diagnostic from the special BPS
mechanism and makes five questions answerable:

\begin{enumerate}
 \item Does a curvature ramp survive when the energy form factor is exactly
 silent?
 \item Does that ramp distinguish random local tangents from structured
 tangents at the same active rank?
 \item Can energy chaos and projector-geometric chaos be varied independently?
 \item Do gauge-invariant response-channel matrix elements approach a
 covariance-matched Wick law along a genuine many-body sequence?
 \item Can non-Abelian holonomy change inside a closed bundle while its
 complete spectrum and Chern class remain fixed?
\end{enumerate}

Our answers establish five complementary advances across spectrum, local
geometry, response tensors, topology, and finite-rank structure.
The metric-normalized curvature is a compression of a fixed signature matrix
by the row space of projector-to-complement response channels.
A Haar row space gives an exact finite-dimensional complex Jacobi point
process, allowing us to derive its connected form factor directly at the
finite rank used in the calculation.
The physical random-local-tangent ensemble follows this prediction over the
registered two-level window.
A complete momentum-resolved tangent family supplies a same-rank structured
reference with only a small number of distinct curvature eigenvalues.  The
adjacent-gap ratio, connected form factor, and number variance place this
reference in a clearly resolved structured regime.

The central evidence is causal rather than correlational.
We vary the Hamiltonian restricted to the degenerate fiber,
\(PHP\), from Poisson to GUE statistics while holding \(P\) and the
projector-to-complement channels fixed; the energy spectrum changes, whereas
the Berry curvature remains exactly invariant.
Separately, we interpolate the tangent channels from Fourier structured to
random local directions at fixed parent Hamiltonian.
Local Jacobi repulsion appears at a smaller interpolation strength than the
full registered form-factor window, while long-range number variance retains
microscopic memory.
The resulting order
\begin{equation}
 \begin{aligned}
 \text{local repulsion}
 &\ \longrightarrow\ \text{two-level ramp}\\
 &\ \longrightarrow\ \text{long-range rigidity}
 \end{aligned}
 \label{eq:headline-hierarchy}
\end{equation}
is a measured hierarchy of geometric correlations.
The form-factor variable is a dimensionless Fourier scale of unfolded
curvature eigenvalues.  Its connection to real-time dynamics defines a
separate and promising extension.

Beyond eigenvalue correlations, fixed local operators expose a resolved
four-channel response cumulant.  Its normalized residual decreases from
\(N=3\) to \(N=5\), demonstrating progressive Gaussianization together with
finite-size operator memory relative to the covariance-matched reference.
On the closed twist torus, a globally periodic ambient conjugation leaves
every energy and the first Chern class invariant while changing Wilson-loop
statistics reproducibly.
The resulting Wilson statistics occupy a structured regime distinct from CUE.
Together, the response tensor and Wilson sector map how local Jacobi behavior
coexists with higher and global microscopic structure.

Finally, an independent rank sequence through \(D=\LargestRankVTwo\) resolves
a finite-dimensional effect unique to signature-compressed ensembles.
When the active rank exceeds the dimension \(M\) of either signature
polarization, exact atoms appear at \(\lambda=\pm1\).
We prove that these deterministic modes reduce the connected plateau from
one to \(k/D\), where \(k=2M-D\) is the continuous-sector dimension, while the
continuous complement retains the Jacobi ramp.
This resolves exact geometric modes and correlated continuous modes within
one curvature matrix.

The paper is organized as follows.
Section~\ref{sec:ensembles} derives the projector response, signature
compression, finite-Jacobi kernel, and atom theorem.
Section~\ref{sec:physical} introduces the fractional topological manifold and
the structured same-rank comparison.
Section~\ref{sec:covariance} performs the independent spectral and geometric
interventions.
Section~\ref{sec:rank} gives the finite-rank atom form factors.
Section~\ref{sec:hierarchy} resolves the correlation hierarchy and its
statistical meaning.
Section~\ref{sec:matrix-topology} tests response-channel Wick factorization
and fixed-Chern Wilson holonomy.
Appendix~\ref{app:numerics} records the numerical and inference protocols.
Appendix~\ref{app:matrix-topology} gives the gauge-covariance, Wick,
Gram-reduction, and bundle-isomorphism derivations.
